% ================================================================
% Appendix: Coderived and contraderived categories
% ================================================================

\section{Coderived and contraderived categories}
\label{app:coderived}
\index{coderived category|textbf}
\index{contraderived category|textbf}

The bar-cobar adjunction at higher genus produces curved complexes
($d^2 \neq 0$), for which the ordinary derived category is trivial.
Positselski's theory of coderived and contraderived
categories~\cite{Positselski11} provides the correct
replacement.  This appendix collects the definitions and key results
needed in this monograph; the chiral specialization appears in
\S\ref{subsec:chiral-coderived-contraderived}.

\subsection{Abstract definitions}
\label{subsec:coderived-abstract}

Let $\mathsf{A}$ be an abelian category with exact direct sums
(for the coderived case) or exact direct products
(for the contraderived case).

\begin{definition}[CDG-algebra and CDG-coalgebra]
\label{def:cdg-algebra-appendix}
\index{CDG-algebra}
\index{CDG-coalgebra}
A \emph{curved differential graded} (CDG) \emph{algebra}
$(A, d, h)$ consists of a graded algebra~$A$, a degree-$1$
map $d \colon A \to A$ satisfying the graded Leibniz rule, and
a degree-$2$ element $h \in A^2$ (the \emph{curvature}) such that
\begin{equation}\label{eq:cdg-relations-appendix}
d^2(a) = [h, a], \qquad d(h) = 0.
\end{equation}
Dually, a \emph{CDG-coalgebra} $(C, d, h)$ consists of a graded
coalgebra~$C$, a degree-$1$ coderivation $d \colon C \to C$,
and a degree-$2$ functional $h \colon C \to k$ such that
$d^2 = h \ast (-)$ and $h \circ d = 0$.
A CDG-algebra (resp.\ coalgebra) with $h = 0$ is an ordinary
DG-algebra (resp.\ DG-coalgebra).
\end{definition}

\begin{definition}[CDG-comodules and CDG-contramodules]
\label{def:cdg-comod-contra-appendix}
\index{CDG-comodule}
\index{CDG-contramodule}
Let $(C, d, h)$ be a CDG-coalgebra.
A \emph{left CDG-comodule} $(M, d_M)$ over~$C$ is a graded left
$C$-comodule equipped with a degree-$1$ map $d_M \colon M \to M$
compatible with the coaction and satisfying $d_M^2 = h \ast (-)$.
A \emph{left CDG-contramodule} $(P, d_P)$ over~$C$ is a graded left
$C$-contramodule with $d_P \colon P \to P$ of degree~$1$ satisfying
the dual curvature relation $d_P^2(p) = h \cdot p$
for all $p \in P$, where $h$ acts via the contraaction.
\end{definition}

\begin{definition}[Coacyclic and contraacyclic objects]
\label{def:coacyclic-contraacyclic}
\index{coacyclic object|textbf}
\index{contraacyclic object|textbf}
Let $(C, d, h)$ be a CDG-coalgebra.
\begin{enumerate}[label=\textup{(\roman*)}]
\item A CDG-comodule $M$ is \emph{coacyclic} if it belongs to the
  minimal thick subcategory of the homotopy category
  $\mathrm{Hot}(C\text{-}\mathrm{comod})$ that contains the
  totalizations of all short exact sequences of CDG-comodules
  and is closed under infinite direct sums.
\item A CDG-contramodule $P$ is \emph{contraacyclic} if it belongs
  to the minimal thick subcategory of
  $\mathrm{Hot}(C\text{-}\mathrm{contra})$ that contains the
  totalizations of all short exact sequences of CDG-contramodules
  and is closed under infinite direct products.
\end{enumerate}
\end{definition}

\begin{definition}[Coderived and contraderived categories;
\ClaimStatusProvedElsewhere]
\label{def:coderived-contraderived-abstract}
\index{coderived category!abstract definition|textbf}
\index{contraderived category!abstract definition|textbf}
The \emph{coderived category} and \emph{contraderived category}
of a CDG-coalgebra $(C, d, h)$ are the Verdier quotients:
\begin{align}
D^{\mathrm{co}}(C\text{-}\mathrm{comod})
  &:= \mathrm{Hot}(C\text{-}\mathrm{comod})
  \,\big/\, \mathrm{Acycl}^{\mathrm{co}}(C\text{-}\mathrm{comod}),
  \label{eq:coderived-abstract}\\
D^{\mathrm{ctr}}(C\text{-}\mathrm{contra})
  &:= \mathrm{Hot}(C\text{-}\mathrm{contra})
  \,\big/\,
  \mathrm{Acycl}^{\mathrm{ctr}}(C\text{-}\mathrm{contra}).
  \label{eq:contraderived-abstract}
\end{align}
These are triangulated categories.  When $h = 0$, every acyclic
object is coacyclic (resp.\ contraacyclic), so there are natural
functors $D^{\mathrm{co}} \to D$ and $D^{\mathrm{ctr}} \to D$
from the exotic to the ordinary derived category.  When
$h \neq 0$, the ordinary derived category is trivial, while
$D^{\mathrm{co}}$ and $D^{\mathrm{ctr}}$ retain nontrivial
structure~\cite[\S\S3--4]{Positselski11}.
\end{definition}

\begin{theorem}[Comodule-contramodule correspondence;
\ClaimStatusProvedElsewhere]
\label{thm:co-contra-correspondence-appendix}
\index{comodule-contramodule correspondence}
Let $(C, d, h)$ be a CDG-coalgebra.  There is an equivalence of
triangulated categories
\begin{equation}\label{eq:co-contra-equivalence}
D^{\mathrm{co}}(C\text{-}\mathrm{comod})
\;\simeq\;
D^{\mathrm{ctr}}(C\text{-}\mathrm{contra}).
\end{equation}
The equivalence is implemented by the functors
$\Phi_C \colon M \mapsto \mathrm{Hom}_C(C, M)$ \textup{(}which
sends injective comodules to projective contramodules\textup{)}
and $\Psi_C \colon P \mapsto C \square_C P$ \textup{(}the
contratensor product\textup{)}.
\end{theorem}

\begin{proof}[Reference]
Positselski~\cite[Theorem~5.3]{Positselski11}.  The chiral
specialization is Theorem~\ref{thm:chiral-co-contra-correspondence}.
\end{proof}

\begin{theorem}[Conilpotent reduction; \ClaimStatusProvedElsewhere]
\label{thm:conilpotent-reduction}
\index{conilpotent coalgebra!coderived reduction}
If the CDG-coalgebra $(C, d, h)$ is \emph{conilpotent}
\textup{(}i.e., the coaugmentation filtration on~$C$ is
exhaustive\textup{)} and $h = 0$, then every acyclic
CDG-comodule is coacyclic, so the coderived category coincides
with the ordinary derived category:
\begin{equation}\label{eq:conilpotent-reduction}
D^{\mathrm{co}}(C\text{-}\mathrm{comod})
\;\simeq\;
D(C\text{-}\mathrm{comod}).
\end{equation}
\end{theorem}

\begin{proof}[Reference]
Positselski~\cite[Corollary~4.4]{Positselski11}.  In this
monograph, the genus-$0$ bar complex
$\bar{B}^{(0)}(\cA)$ is a conilpotent DG-coalgebra, so
the coderived category reduces to the ordinary derived
category---this is the content of
Remark~\textup{\ref{rem:exotic-vs-ordinary}}.
\end{proof}

\begin{remark}[Connection to bar-cobar duality]
\label{rem:coderived-bar-cobar}
\index{bar-cobar adjunction!coderived formulation}
In the setting of this monograph, the bar complex
$\bar{B}^{(g)}(\cA)$ is a chiral CDG-coalgebra with curvature
$m_0^{(g)} = \kappa(\cA) \cdot \omega_g$
(Convention~\ref{conv:higher-genus-differentials}).
At genus~$0$, $m_0^{(0)} = 0$ and
Theorem~\ref{thm:conilpotent-reduction} applies: the coderived
category is the ordinary derived category, and bar-cobar
inversion (Theorem~\ref{thm:higher-genus-inversion})
is a statement in $D(\bar{B}^{(0)}\text{-}\mathrm{comod})$.
At genus $g \geq 1$ with $\kappa(\cA) \neq 0$, the bar complex
is genuinely curved, the ordinary derived category is trivial
(Proposition~\ref{prop:curved-bar-acyclicity}), and bar-cobar
inversion becomes a statement in the coderived category
$D^{\mathrm{co}}(\bar{B}^{(g)}\text{-}\mathrm{comod})$.
The Koszul duality equivalence
$D^{\mathrm{co}}(\cA\text{-}\mathrm{CoMod})
\simeq D^{\mathrm{ctr}}(\cA^!\text{-}\mathrm{ContraMod})$
(Theorem~\ref{thm:positselski-chiral-proved}) is the chiral
incarnation of the comodule-contramodule correspondence
\eqref{eq:co-contra-equivalence}.
\end{remark}

%%% ---------------------------------------------------------------
%%% Relative curved models adequate for the present monograph
%%% ---------------------------------------------------------------

\subsection{Relative curved models}
\label{sec:coderived-models}
\index{coderived category!relative models|textbf}

The full theory of coderived categories of factorization algebras
on $\operatorname{Ran}(X)$ is not yet available in the literature;
its development is part of the Stratum~II programme
(Remark~\ref{rem:two-strata}).  Nevertheless, the off-Koszul
statements in this monograph---where the bar-cobar counit is
a curved equivalence rather than a quasi-isomorphism---require
a precise ambient homotopy category.  This subsection defines the
\emph{minimal} relative homotopy theory needed to make those
statements mathematically honest.

\begin{definition}[Filtered curved factorization model]
\label{def:filtered-curved-model}
A \emph{filtered curved factorization model} is a triple
$(M, d, \mathcal{F})$ where:
\begin{enumerate}[label=\textup{(\roman*)}]
\item $M$ is a conilpotent factorization coalgebra on $\operatorname{Ran}(X)$
  (or a dg module over such);
\item $d \colon M \to M$ is a degree-$1$ map with
  $d^2 = m_0 \in Z(M)$ for a central curvature element~$m_0$
  (cf.\ the fiberwise differential $\dfib$ with $\dfib^{\,2} = \mcurv{g}$,
  Convention~\textup{\ref{conv:higher-genus-differentials}});
\item $\mathcal{F} = \{F^p M\}_{p \geq 0}$ is an exhaustive
  decreasing filtration with $d(F^p) \subset F^p$,
  $m_0 \in F^1$, and $\operatorname{gr}^0_{\mathcal{F}} d$ strict
  ($(\operatorname{gr}^0 d)^2 = 0$).
\end{enumerate}
A morphism $f \colon (M, d, \mathcal{F}) \to (N, d', \mathcal{F}')$ is a
filtration-preserving coalgebra map compatible with differentials and curvatures.
\end{definition}

\begin{definition}[Weak equivalence of curved models]
\label{def:curved-weak-equiv}
A morphism $f \colon M \to N$ of filtered curved factorization models
is a \emph{weak equivalence} if it satisfies either of the following
equivalent conditions:
\begin{enumerate}[label=\textup{(\roman*)}]
\item The associated graded map
  $\operatorname{gr}_{\mathcal{F}} f$ is a quasi-isomorphism of the
  strict complexes $(\operatorname{gr} M, \operatorname{gr}^0 d)
  \xrightarrow{\sim} (\operatorname{gr} N, \operatorname{gr}^0 d')$.
\item The induced map on coacyclic quotients
  $M/\operatorname{coacyc}(M) \to N/\operatorname{coacyc}(N)$ is an
  isomorphism, where $\operatorname{coacyc}(M)$ is the thick subcategory
  generated by totalizations of short exact sequences
  (Positselski~\cite{Positselski11}).
\end{enumerate}
The equivalence of (i) and (ii) holds for bounded-below filtrations
by~\cite[Proposition~3.5]{Positselski11}.
\end{definition}

\begin{definition}[Provisional coderived category]
\label{def:provisional-coderived}
\index{coderived category!provisional|textbf}
The \emph{provisional coderived category} $\mathsf{D}^{\mathrm{co}}_{\mathrm{prov}}(X)$
is the localization of the category of filtered curved factorization
models (Definition~\ref{def:filtered-curved-model}) at the weak
equivalences of Definition~\ref{def:curved-weak-equiv}.
\end{definition}

\begin{proposition}[Adequacy; \ClaimStatusProvedHere]
\label{prop:coderived-adequacy}
The provisional coderived category is adequate for the following
constructions in this monograph:
\begin{enumerate}[label=\textup{(\alph*)}]
\item The curved bar-cobar counit
  $\Omega_X \bar{B}_X(\cA) \to \cA$ of
  Theorem~\textup{\ref{thm:higher-genus-inversion}} is an isomorphism
  in $\mathsf{D}^{\mathrm{co}}_{\mathrm{prov}}(X)$ for every chiral
  algebra~$\cA$ satisfying the Chevalley--Cousin condition.
\item The genus tower $\{\bar{B}^{(g)}(\cA)\}_{g \geq 0}$ of
  Theorem~\textup{\ref{thm:prism-higher-genus}} has compatible
  filtrations making each genus component a filtered curved model.
\item The Verdier involution $\sigma$ acts on filtered curved models
  and descends to $\mathsf{D}^{\mathrm{co}}_{\mathrm{prov}}(X)$,
  so the complementarity decomposition
  \textup{(}Theorem~\textup{\ref{thm:quantum-complementarity-main})}
  holds in the provisional coderived category.
\end{enumerate}
\end{proposition}

\begin{proof}
For (a): the bar-cobar counit is a quasi-isomorphism of associated
graded complexes by the Koszul locus argument of
Theorem~\ref{thm:higher-genus-inversion}, hence a weak equivalence
in the sense of Definition~\ref{def:curved-weak-equiv}(i).
Off the Koszul locus, the counit is still a weak equivalence
in sense~(ii) by the spectral sequence comparison argument of
Theorem~\ref{thm:bar-cobar-spectral-sequence}: the $E_1$ pages agree
and the filtration is bounded below in each bar degree.

For (b): the PBW filtration by conformal weight is exhaustive,
decreasing, preserved by all three components of the bar differential,
and satisfies $m_0^{(g)} \in F^1$ because the curvature involves at
least one OPE contraction.  At associated graded level, the
differential reduces to $\dzero$, which is strict.

For (c): Verdier duality preserves the PBW filtration
(it exchanges weight~$n$ generators with weight~$n$ cogenerators in
the Koszul dual) and commutes with the bar differential.  Hence it
descends to $\mathsf{D}^{\mathrm{co}}_{\mathrm{prov}}(X)$ and the
eigenspace decomposition is well-defined.
\end{proof}

\begin{remark}[Relation to the full theory]
\label{rem:coderived-full}
The provisional coderived category
$\mathsf{D}^{\mathrm{co}}_{\mathrm{prov}}(X)$ is expected to embed
fully faithfully into the eventual coderived category of curved
factorization algebras on $\operatorname{Ran}(X)$.  It is constructed
from the same Positselski localization procedure but restricted to
the class of filtered models appearing in this monograph.  The full
theory would additionally handle pro-nilpotent completions, infinitely
generated bar complexes, and curved $A_\infty$-coalgebras without
filtrations---features needed for the Stratum~II programme but not for
the results proved here.
\end{remark}

\subsection{Completed ambient for factorization algebras}
\label{subsec:completed-ambient}

The full modular programme requires a completed ambient---the
factorization-algebra analogue of the pro-nilpotent/coderived
formalism of Francis--Gaitsgory~\cite{FG12}.

\begin{definition}[Completed chiral ambient; \ClaimStatusConjectured]
\label{def:completed-chiral-ambient}
\index{completed chiral ambient|textbf}
\index{factorization algebra!completed ambient}
Let $X$ be a smooth projective curve.  Define the following full
subcategories of $\Fact(X)$:
\begin{enumerate}[label=\textup{(\roman*)}]
\item $\Alg_{\mathrm{aug,comp}}(\Fact(X))$: the full subcategory of
  augmented factorization algebras \emph{complete} for the
  augmentation-ideal filtration, i.e., those $\cA$ for which the
  canonical map
  $\cA \to \varprojlim_n \cA / \bar{\cA}^{\chirtensor n}$
  is an equivalence.
\item $\CoAlg_{\mathrm{coaug,conil}}(\Fact(X))$: the full subcategory
  of coaugmented factorization coalgebras that are \emph{conilpotent},
  i.e., exhausted by the coradical filtration.
\item $\mathsf{D}^{\mathrm{co}}(\Fact(X))$: the \emph{coderived
  category of curved factorization coalgebras}, defined as the
  localization of the homotopy category at coacyclic objects---those
  generated by totalizations of short exact sequences and closed under
  filtered colimits (cf.\
  Definition~\ref{def:coacyclic-contraacyclic}).
\end{enumerate}
Pro-nilpotence (condition~(i)) is exactly what makes the bar/cobar
differentials converge: the chiral bar construction of an object in
$\Alg_{\mathrm{aug,comp}}$ lands in $\CoAlg_{\mathrm{coaug,conil}}$.
The provisional coderived category
$\mathsf{D}^{\mathrm{co}}_{\mathrm{prov}}(X)$
(Definition~\ref{def:provisional-coderived}) embeds into
$\mathsf{D}^{\mathrm{co}}(\Fact(X))$
as the full subcategory of objects admitting a bounded-below PBW
filtration.
\end{definition}

\begin{remark}[Status]\label{rem:completed-ambient-status}
The $\Eone$-chiral face of this ambient is established: the
monograph proves bar-cobar equivalence in the pro-nilpotent regime
(Theorem~\ref{thm:higher-genus-inversion}).  The full
$\Einf$-factorization version belongs to the Stratum~II programme
(Remark~\ref{rem:two-strata}).  The homotopy form of the modular
Koszul definition
(Definition~\ref{def:modular-koszul-homotopy}) uses
$\Alg_{\mathrm{aug}}$ rather than
$\Alg_{\mathrm{aug,comp}}$ because all objects treated in this
monograph are already complete (their bar complexes converge).
\end{remark}
